
@@M0@@
Le
résultat en découle aussitôt d'après la relation (1.11.3), et l'hypothèse de récurrence.

\textbf{Proposition 5.3.} \uline{Avec les notations et les hypothèses de} 5.1, \uline{on a pour} \uline{tout} @@M1@@ \uline{une décomposition en somme directe de} @@M2@@-\uline{modules} :
@@M3@@
\uline{la filtration considérée étant la filtration de} @@M4@@-\uline{algèbre augmentée définie} \uline{en} V 3.9.1 \uline{et} 3.10.

On remarque d'abord que @@M5@@ est d'augmentation nulle. Par ailleurs, on a :
@@M6@@
Donc @@M7@@ pour tout @@M8@@. Par suite, @@M9@@ est un quotient de la @@M10@@-@@M11@@-algèbre engendrée par un générateur @@M12@@ soumis aux relations @@M13@@ pour tout @@M14@@, et l'idéal par lequel on fait le quotient est engendré par la relation (1.13.1). On peut donc appliquer V 4.12 et V 4.13.ii), et on est ramené à montrer que pour tout @@M15@@ :
@@M16@@

C'est évident pour @@M17@@; on poursuit par récurrence. La relation (1.13.1) donne :
@@M18@@
appliquant l'hypothèse de récurrence aux @@M19@@, on en déduit le résultat cherché, la somme étant directe puisque @@M20@@ sont indépendants sur @@M21@@.

